% === Method ===
\section{Method}
\label{sec:method}

In this section, we describe our proposed method. We first formulate the problem in \cref{sec:formulation}, then present the overall architecture in \cref{sec:architecture}, and finally describe the training procedure in \cref{sec:training}.

% --- Example: how to include a figure ---
\begin{figure}[t]
  \centering
  \fbox{\rule{0pt}{2in}\rule{0.9\linewidth}{0pt}}  % placeholder box; replace with:
  % \includegraphics[width=\linewidth]{figures/your_figure.pdf}
  \caption{Overview of the proposed \method framework. Replace this placeholder with your method pipeline figure.}
  \label{fig:method_overview}
\end{figure}

% --- Example: how to write an equation ---
\subsection{Problem Formulation}
\label{sec:formulation}

Given a dataset $\mathcal{D} = \{(x_i, y_i)\}_{i=1}^{N}$, where $x_i$ denotes the input and $y_i$ denotes the label, our goal is to learn a mapping $f: \mathcal{X} \rightarrow \mathcal{Y}$ that minimizes the expected loss:
\begin{equation}
  \mathcal{L} = \mathbb{E}_{(x,y) \sim \mathcal{D}} \left[ \ell\bigl(f(x; \theta),\, y\bigr) \right]
  \label{eq:objective}
\end{equation}
where $\theta$ denotes the model parameters and $\ell(\cdot, \cdot)$ is the loss function. We refer to \cref{eq:objective} as the training objective throughout this paper.

% --- Example: multi-line equation with a SINGLE number ---
% Use "aligned" inside "equation" for one number; use "align" for per-line numbers.
\begin{equation}
  \begin{aligned}
    \nabla_\theta \mathcal{L}
    &= \frac{1}{N} \sum_{i=1}^{N} \nabla_\theta \ell\bigl(f(x_i; \theta),\, y_i\bigr) \\
    &\approx \frac{1}{|\mathcal{B}|} \sum_{i \in \mathcal{B}} \nabla_\theta \ell\bigl(f(x_i; \theta),\, y_i\bigr),
  \end{aligned}
  \label{eq:gradient}
\end{equation}
where $\mathcal{B}$ denotes a mini-batch sampled from $\mathcal{D}$.

\subsection{Architecture}
\label{sec:architecture}

Describe the architecture of your model here. Explain each component and how they interact with each other.

% --- Example: pseudocode ---
\begin{algorithm}[t]
  \caption{Training Procedure of \method}
  \label{alg:training}
  \KwIn{Dataset $\mathcal{D}$, learning rate $\eta$, number of epochs $T$}
  \KwOut{Trained model parameters $\theta$}
  Initialize $\theta$ randomly\;
  \For{$t = 1$ \KwTo $T$}{
    Sample mini-batch $\mathcal{B} \subset \mathcal{D}$\;
    Compute loss $\mathcal{L}(\theta; \mathcal{B})$ via \cref{eq:objective}\;
    Compute gradient $g \leftarrow \nabla_\theta \mathcal{L}$\;
    Update $\theta \leftarrow \theta - \eta \cdot g$\;
    \If{convergence criterion met}{
      \textbf{break}\;
    }
  }
  \Return{$\theta$}
\end{algorithm}

The training procedure of \method is summarized in \cref{alg:training}.

\subsection{Training Procedure}
\label{sec:training}

Describe how the model is trained, including the optimizer, learning rate schedule, data augmentation, and any other relevant details.
